\documentclass[12pt,a4paper]{report}

% Packages utiles
\usepackage[utf8]{inputenc}
\usepackage[T1]{fontenc}
\usepackage[french]{babel}
\usepackage{lmodern}
\usepackage{geometry}
\usepackage{setspace}
\usepackage{hyperref}

\geometry{margin=2.5cm}
\setlength{\parskip}{0.5em}
\setlength{\parindent}{0pt}
\onehalfspacing

% Page de garde
\title{\Huge Rapport 2 \\[0.5cm]
\Large Systèmes d'Exploitation – NACHOS \\[0.2cm]
\large Devoir 2 : Gestion des threads utilisateurs}
\author{Mervyn Samsan - Robert Dylan - Negre Damien}
\date{\today}

\begin{document}

\maketitle
\newpage

\section*{I. Bilan}
\addcontentsline{toc}{section}{I. Bilan}

- L’objectif de ce TP était d’implémenter la gestion des threads utilisateurs dans NachOS, à travers les appels systèmes \texttt{ThreadCreate} et \texttt{ThreadExit}. \\
- Nous avons réussi à faire fonctionner ces appels systèmes, aussi bien côté noyau que côté utilisateur. \\
- Les threads peuvent désormais s’exécuter en parallèle dans le même espace mémoire, sans interférer entre eux. \\
- Chaque thread dispose de sa propre pile grâce aux fonctions \texttt{AllocateUserStack} et \texttt{FreeUserStack}, ce qui évite les conflits entre variables locales. \\
- La terminaison du processus est correctement gérée : le programme ne s’arrête que lorsque le dernier thread actif appelle \texttt{ThreadExit}, grâce à un compteur de threads dans l’\texttt{AddrSpace}. \\
- Nous avons également implémenté la terminaison automatique en plaçant l’adresse de \texttt{ThreadExit} dans le registre \$ra lors de la création du thread, ce qui assure un arrêt propre. \\
-Enfin, nous avons également implémenté la dernière partie concernant les sémaphores utilisateurs. Notre test producteur-consommateur est fonctionnel et démontre une synchronisation correcte, s'exécutant en boucle continue et affichant la succession attendue de 'P' (Producteur) et 'C' (Consommateur). \\

\section*{II. Points délicats}
\addcontentsline{toc}{section}{II. Points délicats}

- Ce TP a été plus complexe que le premier, surtout pour bien comprendre la gestion du contexte d’exécution et de la mémoire dans NachOS. \\
- La fonction \texttt{StartUserThread} a été la partie la plus difficile à mettre en place. Il fallait créer un nouveau contexte pour le thread tout en gardant le même espace mémoire que le thread principal. \\
- Le passage des paramètres a également été un point sensible. Comme \texttt{Thread::Start} ne prend qu’un seul argument, nous avons créé une structure \texttt{UserThreadParams} regroupant plusieurs informations : l’adresse de la fonction à exécuter, son argument, l’adresse de sortie (\texttt{exitAddr}) et le sommet de la pile (\texttt{stackTop}). Cette structure est allouée dans \texttt{do\_ThreadCreate} puis libérée au début de \texttt{StartUserThread}. \\
- Il fallait ensuite :
    \begin{itemize}
        \item Allouer une pile spécifique pour le thread via \texttt{AllocateUserStack()} ;
        \item Initialiser tous les registres MIPS à zéro pour éviter les comportements imprévisibles ;
        \item Configurer correctement les registres importants : le \texttt{PCReg}, le \texttt{NextPCReg}, le \texttt{StackReg}, le registre d’argument (\$a0) et le registre d’adresse de retour (\$ra). \\
    \end{itemize}
- La moindre erreur dans ces registres provoquait un crash du programme, ce qui a rendu le débogage assez long. \\
 
Pour corriger les erreurs, nous avons d’abord ajouté des \texttt{printf} et des \texttt{PutChar} pour suivre le déroulement du code.
Notre chargé de TD nous a ensuite expliqué la logique à adopter pour localiser précisément l’origine des erreurs, en nous montrant notamment comment analyser le déroulement du contexte et des registres.
Nous avons également utilisé les options de débogage intégrées à NachOS (notamment \texttt{-d} et \texttt{-d t}), qui se sont révélées très utiles, non seulement à ce moment-là, mais aussi plus tard dans le projet pour comprendre et résoudre d’autres problèmes.


\section*{III. Limitations}
\addcontentsline{toc}{section}{III. Limitations}

- Notre implémentation fonctionne correctement, mais présente encore quelques limites malgré l’ajout des sémaphores utilisateurs pour la synchronisation. \\
- Cela implique que si plusieurs threads accèdent à une même variable en même temps, des conditions de course peuvent apparaître. \\
- Le système devient également instable si trop de threads sont créés, car la taille du \texttt{Bitmap} gérant les piles est fixe. \\

\section*{IV. Tests}
\addcontentsline{toc}{section}{IV. Tests}

Tous nos tests ont été réalisés dans le fichier \texttt{makethreads.c}, que nous avons utilisé pour valider chaque étape de notre implémentation. \\

- \textbf{Test de base :} Nous avons commencé par créer un seul thread affichant un caractère pendant que le main restait dans une boucle infinie. Cela nous a permis de vérifier le bon fonctionnement de \texttt{ThreadCreate} et \texttt{StartUserThread}.\\
- \textbf{Test de concurrence :} Nous avons ensuite lancé plusieurs threads affichant des messages différents (par exemple “AAAAA” et “BBBBB”). Le mélange des sorties à l’écran nous a confirmé que les threads s’exécutaient bien en parallèle.\\
- \textbf{Test de stress :} Nous avons créé 3 threads, chacun affichant deux lettres ’a’. Nous avons bien obtenu au total 6 caractères ’a’ à l’écran, ce qui prouve que tous les threads se sont exécutés correctement et indépendamment. Aucun crash n’a été observé, ce qui confirme la bonne gestion des piles et du partage de mémoire. \\
- \textbf{Test de terminaison :} Enfin, nous avons vérifié que le programme ne se termine qu’après la fin du dernier thread actif. Après que tous les threads enfants aient appelé \texttt{ThreadExit}, le processus principal s’arrêtait proprement.\\

Ces différents tests nous ont permis de valider le bon fonctionnement global de la gestion des threads utilisateurs dans NachOS. \\


\end{document}
